\documentclass[11pt,a4paper]{article}
\usepackage[T1]{fontenc}
\usepackage[utf8]{inputenc}
\usepackage{lmodern}
\usepackage{amsmath,amssymb,amsthm}
\usepackage{booktabs}
\usepackage{xcolor}
\usepackage[margin=27mm]{geometry}
\usepackage[colorlinks=true,urlcolor=blue,citecolor=blue,linkcolor=blue]{hyperref}
\newtheorem{theorem}{Theorem}
\newtheorem{proposition}[theorem]{Proposition}
\newtheorem{lemma}[theorem]{Lemma}
\newtheorem{corollary}[theorem]{Corollary}
\theoremstyle{definition}
\newtheorem{definition}[theorem]{Definition}
\newtheorem{remark}[theorem]{Remark}
\theoremstyle{plain}
\DeclareMathOperator{\Tr}{Tr}
\newcommand{\ket}[1]{\lvert#1\rangle}
\newcommand{\bra}[1]{\langle#1\rvert}
\newcommand{\norm}[1]{\lVert#1\rVert}
\newcommand{\one}{\mathbf{1}}
\newcommand{\cL}{\mathcal{L}}
\newcommand{\cD}{\mathcal{D}}
\setlength{\emergencystretch}{2em}
\setlength{\parindent}{0pt}
\setlength{\parskip}{0.3em plus 0.1em minus 0.1em}
\title{A Mpemba mode invisible to every strict reduction}
\author{Fable \& Astra\\[0.35em]
\normalsize Research led by Jean-Philippe Bruneton\\[-0.1em]
\small Universit\'e Paris Cit\'e, France}
\date{19 November 2026}

\begin{document}
\maketitle
\begin{abstract}
In the Mpemba effect a system that starts farther from equilibrium reaches
it sooner than one that starts closer. In Markovian open quantum systems the
effect is governed by the slowest relaxation mode of the generator: a state
with a small component along that mode can overtake a state that starts
closer to equilibrium but carries a larger slow component. Recent work
proposes to witness the effect through the expectation value of a single
observable. We ask how large the support of such an observable must be in a
many-body system, and show that it can be forced to be the whole system. We
give an explicit Lindbladian on a chain of $n\ge2$ qubits, with jump
operators acting on one site or on two neighbouring sites, whose stationary
state is unique and whose slowest mode is the global parity
$Z_1\cdots Z_n$, at rate $2\epsilon$ with $0<\epsilon<1$. The full spectrum
is $\{-2(m+\epsilon b)\}$ with multiplicity $\binom{2n-1}{m}$, for
$0\le m\le 2n-1$ and $b\in\{0,1\}$. The state $(\one+\alpha Z_1\cdots
Z_n)/2^n$ has maximally mixed reductions on every strict subset of sites at
all times, yet it relaxes at the slowest rate, and it is overtaken by the
state $(\one+\beta X_n)/2^n$ with $\beta>\alpha$ at time
$\log(\beta/\alpha)/(2-2\epsilon)$, in trace distance, in relative entropy
and in every unitarily invariant distance to equilibrium that grows with the
deviation from the maximally mixed state. Any observable sensitive
to the slow amplitude acts on all $n$ sites.
\end{abstract}

\section{Introduction}
A cup of hot water left outside can freeze before a cup of warm water. This
counterintuitive ordering, the Mpemba effect, has a clean formulation for
Markovian dynamics \cite{LuRaz2017}. Relaxation towards a unique stationary
state is a superposition of decaying modes, and at late times only the
slowest one matters. Two initial states are then ranked by the size of their
component along the slowest mode, not by their initial distance to
equilibrium. A state farther from equilibrium but with a smaller slow
component eventually becomes the closer one, and the two distance curves
cross. In open quantum systems, Carollo, Lasanta and Lesanovsky
\cite{Carollo2021} used this mechanism to accelerate relaxation
exponentially, by applying a unitary rotation that removes the component of
the initial state along the slowest mode.

Detecting the effect requires access to the slow component. Choquehuanca,
Obando, Sarandy and de Paula \cite{Choquehuanca2026} proposed to witness the
Mpemba crossing with the expectation value of a single Hermitian observable,
whose expansion on the eigenmodes of the generator is tuned to the slow and
fast modes, and illustrated it on a single qubit. In a many-body system the
natural follow-up question concerns the support of such an observable. If
the slowest mode were visible in the state of one site or of a few
neighbouring sites, a local probe would suffice. This note asks whether a
local and primitive dynamics on a chain can have a slow mode that is global
in the strongest sense, that is, invisible in the reduced state of every
strict subset of sites.

The answer is yes, and the construction is elementary. The dynamics on $n$
qubits combines nearest-neighbour bit flips $X_jX_{j+1}$, single-site
dephasing $Z_j$, and a weak bit flip $X_1$ at rate $\epsilon$ on the first
site. Without the boundary term the global parity $P=Z_1\cdots Z_n$ commutes
with every jump operator and would be conserved; the weak boundary term
breaks this symmetry and turns $P$ into the unique slowest mode, at rate
$2\epsilon$. Since $P$ is a product of traceless factors on all sites, every
partial trace over a nonempty set of sites annihilates it. The state
$\rho_P=(\one+\alpha P)/2^n$ is therefore indistinguishable from equilibrium
on every strict subsystem at all times, while it is the slowest-relaxing
state of the chain. The state $\rho_Q=(\one+\beta X_n)/2^n$, with
$\beta>\alpha$, starts farther from equilibrium, visibly so on the last
site, relaxes at rate $2$, and overtakes $\rho_P$ at
$t^*=\log(\beta/\alpha)/(2-2\epsilon)$.

The main results are the exact spectrum of the generator
(Proposition~\ref{prop:spectrum}), the uniqueness of the stationary state
and of the slowest mode (Corollary~\ref{cor:slow}), the invisibility of the
slow mode in strict reductions (Proposition~\ref{prop:invisible}), and the
Mpemba crossing with its time (Theorem~\ref{thm:crossing}). The example
separates two notions that are easily conflated. The slow amplitude can be
estimated by measuring every qubit locally in the $Z$ basis and multiplying
the outcomes, so local measurements suffice once their records are
correlated; what fails is every observable, and every reduced state, whose
support is a strict subset of the chain.

\section{Related work}\label{sec:related}
Lu and Raz \cite{LuRaz2017} formulated the Mpemba effect and its inverse
for Markovian dynamics and derived a sufficient condition for it within
nonequilibrium thermodynamics. We work in the same Markovian setting, where
late-time relaxation is controlled by the slowest eigenmode of the generator.
Carollo, Lasanta and Lesanovsky \cite{Carollo2021} brought the effect to
Markovian open quantum systems, where the generator has the
Gorini--Kossakowski--Sudarshan--Lindblad form, and
showed that a unitary rotation making the initial state orthogonal to the
slowest mode yields an exponential speed-up, with applications to Dicke
models and all-to-all interacting spins. Such a rotation is cheap in the
present example: a Hadamard gate on site $1$ maps $P$ to the faster string
$X_1Z_2\cdots Z_n$. What is global here is the detection of the slow
component, not its removal.

Choquehuanca, Obando, Sarandy and de Paula \cite{Choquehuanca2026} witness
the quantum Mpemba effect through a single observable whose eigenmode
coefficients are tuned to the slow and fast modes, and they demonstrate the
construction on a qubit under generalized amplitude damping; many-body
extensions are left to future work. Their construction uses a nonzero
coefficient on the slowest mode. The construction below shows that in a
local many-body dynamics this coefficient can force the observable to act on
every site (Remark~\ref{rem:witness}), and it gives a fully explicit
many-body test case in which the slow and fast modes and the crossing time
are known in closed form.

\section{The dynamics and its spectrum}\label{sec:spectrum}
For a Hermitian unitary $U$ on $(\mathbb C^2)^{\otimes n}$ let
$\cD_U(\rho)=U\rho U^\dagger-\rho$, the dissipator of a single jump operator
$U$ at unit rate, since $U^\dagger U=\one$. For $n\ge2$ and $0<\epsilon<1$
consider
\begin{equation}\label{eq:L}
\cL=\sum_{j=1}^{n-1}\cD_{X_jX_{j+1}}+\epsilon\,\cD_{X_1}
+\sum_{j=1}^n\cD_{Z_j},
\end{equation}
where $X_j,Y_j,Z_j$ are the Pauli matrices on site $j$. Every jump operator
acts on one site or on two neighbouring sites, so $\cL$ is local. A Pauli
string is a product $\sigma=\sigma_1\otimes\cdots\otimes\sigma_n$ with
$\sigma_j\in\{I,X,Y,Z\}$. For such a string define its dephasing weight
$k(\sigma)=\#\{j:\sigma_j\in\{X,Y\}\}$, its flip pattern
$f_j(\sigma)=1$ if $\sigma_j\in\{Y,Z\}$ and $0$ otherwise, and its number of
domain walls $e(\sigma)=\#\{j\le n-1: f_j(\sigma)\ne f_{j+1}(\sigma)\}$.

\begin{proposition}\label{prop:spectrum}
Every Pauli string is an eigenvector of $\cL$:
\[
\cL(\sigma)=-r(\sigma)\,\sigma,\qquad
r(\sigma)=2k(\sigma)+2e(\sigma)+2\epsilon f_1(\sigma).
\]
Consequently $\cL$ is diagonalizable with real spectrum consisting of the
numbers $-2(m+\epsilon b)$, $0\le m\le 2n-1$, $b\in\{0,1\}$, the eigenvalue
$-2(m+\epsilon b)$ having multiplicity $\binom{2n-1}{m}$.
\end{proposition}

\begin{proof}
Two Pauli strings either commute or anticommute. If $U$ is a Pauli string,
$U\sigma U^\dagger=\sigma$ when they commute and $-\sigma$ when they
anticommute, so $\cD_U(\sigma)$ equals $0$ or $-2\sigma$. The jump $Z_j$
anticommutes with $\sigma$ exactly when $\sigma_j\in\{X,Y\}$, which
contributes $2k(\sigma)$. The jump $X_1$ anticommutes with $\sigma$ exactly
when $f_1(\sigma)=1$, which contributes $2\epsilon f_1(\sigma)$. The jump
$X_jX_{j+1}$ anticommutes with $\sigma$ exactly when exactly one of
$\sigma_j,\sigma_{j+1}$ anticommutes with $X$, that is, when
$f_j(\sigma)\ne f_{j+1}(\sigma)$, which contributes $2e(\sigma)$. The $4^n$
Pauli strings form a basis of the operator space, so $\cL$ is diagonal in
that basis.

For the multiplicities, a Pauli string is determined by its flip pattern
$f\in\{0,1\}^n$ together with the choice, on each site, between the two
letters with that flip value, one of which is dephased ($X$ or $Y$) and one
not ($I$ or $Z$). The dephasing weight is therefore independent of the flip
pattern, and $k$ takes the value $k$ on $\binom nk$ choices. A flip pattern
is determined by $f_1$ and the positions of its domain walls among the $n-1$
bonds, so for each value $b$ of $f_1$ there are $\binom{n-1}{e}$ patterns
with $e$ walls. The number of strings with $r=2(m+\epsilon b)$ is thus
$\sum_{k+e=m}\binom nk\binom{n-1}{e}=\binom{2n-1}{m}$ by Vandermonde's
identity. Since $\epsilon\notin\mathbb Z$, distinct pairs $(m,b)$ give
distinct eigenvalues.
\end{proof}

\begin{corollary}\label{cor:slow}
The identity spans the kernel of $\cL$, so $\one/2^n$ is the unique
stationary state, and every other eigenvalue has strictly negative real
part; the dynamics is primitive. The spectral gap is $2\epsilon$, the
eigenvalue $-2\epsilon$ is simple with eigenvector $P=Z_1Z_2\cdots Z_n$, and
every other nonzero eigenvalue is at most $-2$. In particular the slowest
non-stationary mode is unique and global.
\end{corollary}

\begin{proof}
By Proposition~\ref{prop:spectrum} the eigenvalues $0$ and $-2\epsilon$
correspond to $(m,b)=(0,0)$ and $(0,1)$, each with multiplicity
$\binom{2n-1}{0}=1$; all others have $m\ge1$ and are at most $-2<-2\epsilon$.
A string with $m=0$ has $k=0$ and $e=0$, so it uses only $I$ and $Z$ with a
constant flip pattern: it is $\one$ when $f_1=0$ and $P$ when $f_1=1$.
\end{proof}

The role of the boundary term is now clear. For $\epsilon=0$ the parity $P$
commutes with every jump operator and the stationary space is spanned by
$\one$ and $P$. The term $\epsilon\cD_{X_1}$ breaks this conservation law
weakly and leaves $P$ as a long-lived mode; the condition $\epsilon<1$ makes
it slower than every other mode, the next rate being $2$, attained for
instance by $X_n$, which has $k=1$, $e=0$ and $f_1=0$.

\section{Invisibility in strict reductions}\label{sec:invisible}
For a subset $A\subsetneq\{1,\dots,n\}$ write $\Tr_{\bar A}$ for the partial
trace over the complementary sites $\bar A\ne\emptyset$, and expand an
operator as $\rho=2^{-n}\sum_\sigma c_\sigma(\rho)\,\sigma$ with
$c_\sigma(\rho)=\Tr(\sigma\rho)$. By Corollary~\ref{cor:slow}, the amplitude
of the slowest mode in $\rho$ is $c_P(\rho)=\Tr(P\rho)$, and under the
dynamics $c_P(e^{t\cL}\rho)=e^{-2\epsilon t}c_P(\rho)$.

\begin{proposition}\label{prop:invisible}
For every strict subset $A$ of sites, $\Tr_{\bar A}P=0$. Consequently:
\begin{enumerate}
\item for every state $\rho$ and every real $\delta$ such that
$\rho+\delta P/2^n\ge0$, the states $\rho$ and $\rho+\delta P/2^n$ have the
same reductions to every strict subset of sites but slow amplitudes
differing by $\delta$;
\item for $\rho_P=(\one+\alpha P)/2^n$ with $|\alpha|\le1$, the evolved
state $e^{t\cL}\rho_P=(\one+\alpha e^{-2\epsilon t}P)/2^n$ has maximally
mixed reductions $\one_A/2^{|A|}$ to every strict subset $A$, at all times
$t\ge0$;
\item an observable $M$ satisfies $\Tr(MP)\ne0$, that is, its expectation
value depends on the slow amplitude, only if $M$ acts nontrivially on every
site.
\end{enumerate}
\end{proposition}

\begin{proof}
The partial trace of a product operator factorizes, and $\Tr Z_j=0$ for
every traced site, so $\Tr_{\bar A}P=0$ whenever $\bar A\ne\emptyset$. Items
(1) and (2) follow by linearity, using Corollary~\ref{cor:slow} for the time
dependence. For (3), if $M$ acts as the identity on some site $j$, then
$M=M'\otimes\one_j$ and $\Tr(MP)=\Tr(M'\,\Tr_j P)=0$.
\end{proof}

\section{The Mpemba crossing}\label{sec:crossing}
Let $0<\alpha<\beta<1$ and
\[
\rho_P=\frac{\one+\alpha P}{2^n},\qquad
\rho_Q=\frac{\one+\beta X_n}{2^n}.
\]
Both are states: $P$ and $X_n$ are Hermitian with eigenvalues $\pm1$, so the
eigenvalues of $\rho_P$ and $\rho_Q$ are $(1\pm\alpha)/2^n$ and
$(1\pm\beta)/2^n$, each with multiplicity $2^{n-1}$. Write
$\rho_P(t)=e^{t\cL}\rho_P$, $\rho_Q(t)=e^{t\cL}\rho_Q$, and let
$T(\rho)=\frac12\norm{\rho-\one/2^n}_1$ denote the trace distance to
equilibrium.

\begin{theorem}\label{thm:crossing}
For all $t\ge0$,
\[
T(\rho_P(t))=\tfrac12\alpha e^{-2\epsilon t},\qquad
T(\rho_Q(t))=\tfrac12\beta e^{-2t}.
\]
The state $\rho_Q$ starts farther from equilibrium and is strictly closer
to it for all $t>t^*$, where
\[
t^*=\frac{\log(\beta/\alpha)}{2-2\epsilon}.
\]
The same crossing, at the same time $t^*$ and with the same reversal of the
order afterwards, holds for every distance to $\one/2^n$ of the form
$h(\rho)$ with $h$ unitarily invariant and strictly increasing along
$a\mapsto(\one+aU)/2^n$, $a\in[0,1)$, for a Pauli string $U\ne\one$; this
includes the Hilbert--Schmidt distance and the relative entropy
$S(\rho\,\|\,\one/2^n)$. Throughout, every strict reduction of
$\rho_P(t)$ is maximally mixed, whereas the reduction of $\rho_Q(t)$ to
site $n$ is $(\one+\beta e^{-2t}X)/2$.
\end{theorem}

\begin{proof}
By Proposition~\ref{prop:spectrum}, $\cL(P)=-2\epsilon P$ and
$\cL(X_n)=-2X_n$ (for $X_n$: $k=1$, $e=0$, $f_1=0$ since $n\ge2$), and
$\cL(\one)=0$. Hence $\rho_P(t)=(\one+a_P(t)P)/2^n$ and
$\rho_Q(t)=(\one+a_Q(t)X_n)/2^n$ with $a_P(t)=\alpha e^{-2\epsilon t}$ and
$a_Q(t)=\beta e^{-2t}$. For a Hermitian unitary $U\ne\one$ that is a Pauli
string, $\norm{aU}_1=|a|\,2^n$, which gives the trace distances. The
difference $\log a_Q(t)-\log a_P(t)=\log(\beta/\alpha)-(2-2\epsilon)t$ is
strictly decreasing, positive at $t=0$, and vanishes at $t^*$, so
$a_Q(t)<a_P(t)$ exactly for $t>t^*$. Any two Pauli strings different from
$\one$ are unitarily conjugate, so $\rho_P(t)$ and $\rho_Q(t)$ are unitarily
conjugate states of the family $(\one+aU)/2^n$ evaluated at $a_P(t)$ and
$a_Q(t)$; a unitarily invariant $h$ that is strictly increasing in $a$
therefore orders them exactly as $a_P$ and $a_Q$ do. For the Hilbert--Schmidt
distance, $\norm{aU/2^n}_2=a\,2^{-n/2}$. For the relative entropy,
$S\big((\one+aU)/2^n\,\big\|\,\one/2^n\big)
=\frac12\big[(1+a)\log(1+a)+(1-a)\log(1-a)\big]$, whose derivative
$\frac12\log\frac{1+a}{1-a}$ is positive for $0<a<1$. The statements on
reductions are Proposition~\ref{prop:invisible} for $\rho_P$ and a direct
partial trace for $\rho_Q$.
\end{proof}

For example, with $\epsilon=1/6$, $\alpha=0.3$ and $\beta=0.8$, the
crossing time is $t^*=\tfrac35\log\tfrac83\approx0.5885$.

The theorem makes the separation quantitative. A monitor that sees only
strict subsystems declares $\rho_P(t)$ at equilibrium for all $t$, and sees
$\rho_Q(t)$ out of equilibrium on the last site; it orders the two states
opposite to their late-time global order and misses the existence of the
slow tail altogether. The slow tail of $\rho_P$ is real nonetheless: its
trace distance to equilibrium decays only as $e^{-2\epsilon t}$, and for
$\epsilon\to0$ its relaxation time $1/(2\epsilon)$ diverges.

\begin{remark}[single-observable witnesses]\label{rem:witness}
In the framework of \cite{Choquehuanca2026}, an observable $M$ witnesses the
crossing through its coefficients on the slow and fast eigenmodes, and the
construction uses a nonzero coefficient on the slow mode. For the generator
\eqref{eq:L} the slow mode is $P$, so the witness needs $\Tr(MP)\ne0$, and by
Proposition~\ref{prop:invisible}(3) it acts nontrivially on all $n$ sites.
This is a statement about supports, not about the difficulty of the
measurement: $\Tr(P\rho)$ is the mean of the product of the outcomes of
single-site $Z$ measurements on one copy of $\rho$, so local measurements
with correlated classical post-processing estimate it.
\end{remark}

\section{Verification}\label{sec:verification}
The
\href{https://github.com/jpbruneton/open_quantum_problems/blob/main/research/entries/locally-invisible-mpemba-mode/checks/verify_mpemba_mode.py}{Python checker}
assembles $\cL$ as a $4^n\times4^n$ matrix for $n=2,3,4,5$ and
$\epsilon\in\{1/6,1/2,9/10\}$, diagonalizes it numerically, and compares
the result, with multiplicities, to the rate formula and to the closed form
of Proposition~\ref{prop:spectrum}, which it evaluates in rational
arithmetic. It checks Corollary~\ref{cor:slow} exactly from the rates. For
$n\le4$ it evolves $\rho_P$ and $\rho_Q$ with the matrix exponential of
$\cL$, checks the trace distances of Theorem~\ref{thm:crossing}, locates the
crossing by root finding and compares it with $t^*$, checks that
$\rho_P(t^*)$ and $\rho_Q(t^*)$ are isospectral, and checks that every
strict reduction of $\rho_P(t)$ is maximally mixed. Its
\href{https://github.com/jpbruneton/open_quantum_problems/blob/main/research/entries/locally-invisible-mpemba-mode/data/verification.json}{recorded output}
lists the exact spectra. All deviations are below $10^{-9}$, and the checker
runs in about ten seconds.

\section{Scope and limitations}\label{sec:scope}
The proofs are analytic and cover every $n\ge2$ and every $0<\epsilon<1$.
The numerical checks confirm the conventions and the closed forms for
$n\le5$; they are not a formal verification.

The obstruction concerns observables and reduced states supported on strict
subsets of sites, measured on a single copy. It does not say that the slow
amplitude is hard to measure: products of local outcomes give it directly,
as noted in Remark~\ref{rem:witness}. The construction is a single explicit
family; it does not characterize which local Lindbladians have a slowest
mode invisible to strict reductions, nor does it address thermal (non-unital)
dynamics, where the stationary state is not maximally mixed.

On bibliographic priority, the note relies on the works cited above and no
systematic search for antecedents was documented. The individual
ingredients, Pauli-diagonal dissipators and global parity as a weakly broken
conservation law, are standard, and we make no claim of priority for the
construction.

\begin{thebibliography}{9}
\bibitem{LuRaz2017}
Z.~Lu and O.~Raz,
\emph{Nonequilibrium thermodynamics of the Markovian Mpemba effect and its
inverse},
Proc. Natl. Acad. Sci. U.S.A. \textbf{114}, 5083--5088 (2017),
\href{https://doi.org/10.1073/pnas.1701264114}{doi:10.1073/pnas.1701264114},
\href{https://arxiv.org/abs/1609.05271}{arXiv:1609.05271}.

\bibitem{Carollo2021}
F.~Carollo, A.~Lasanta and I.~Lesanovsky,
\emph{Exponentially accelerated approach to stationarity in Markovian open
quantum systems through the Mpemba effect},
Phys. Rev. Lett. \textbf{127}, 060401 (2021),
\href{https://doi.org/10.1103/PhysRevLett.127.060401}{doi:10.1103/PhysRevLett.127.060401},
\href{https://arxiv.org/abs/2103.05020}{arXiv:2103.05020}.

\bibitem{Choquehuanca2026}
J.~M.~Z. Choquehuanca, P.~A.~C. Obando, M.~S. Sarandy and F.~M. de~Paula,
\emph{Witnessing the quantum Mpemba effect with a single observable},
preprint (2026),
\href{https://arxiv.org/abs/2607.12163}{arXiv:2607.12163}.

\end{thebibliography}
\end{document}
